\documentclass[11pt]{article}

\usepackage[margin=1in]{geometry}
\usepackage[T1]{fontenc}
\usepackage[utf8]{inputenc}
\usepackage{lmodern}
\usepackage{microtype}
\usepackage[table]{xcolor}
\usepackage{amsmath,amssymb,mathtools,bm}
\usepackage{booktabs,array,longtable,tabularx}
\usepackage{hyperref}
\usepackage{enumitem}
\usepackage{fancyhdr}
\usepackage{titlesec}
\usepackage{tcolorbox}
\usepackage{ragged2e}
\usepackage{setspace}
\hypersetup{colorlinks=true,linkcolor=blue!45!black,urlcolor=blue!45!black,citecolor=blue!45!black}
\definecolor{TitleBlue}{HTML}{163A5F}
\definecolor{AccentTeal}{HTML}{2D6F73}
\definecolor{SoftBlue}{HTML}{EAF3F8}
\definecolor{SoftTeal}{HTML}{EDF8F6}
\definecolor{SoftGray}{HTML}{F5F7FA}
\pagestyle{fancy}
\fancyhf{}
\lhead{PINN\_M3DC1}
\rhead{Project Plan}
\cfoot{\thepage}
\fancypagestyle{plain}{%
  \fancyhf{}%
  \lhead{PINN\_M3DC1}%
  \rhead{Project Plan}%
  \cfoot{\thepage}%
}
\setlength{\headheight}{14pt}
\setlength{\parskip}{0.45em}
\setlength{\parindent}{0pt}
\onehalfspacing
\numberwithin{equation}{section}
\titleformat{\section}{\Large\bfseries\color{TitleBlue}}{\thesection}{0.6em}{}
\titleformat{\subsection}{\large\bfseries\color{AccentTeal}}{\thesubsection}{0.6em}{}
\newtcolorbox{overviewbox}{colback=SoftBlue,colframe=TitleBlue,boxrule=0.8pt,arc=2mm,left=3mm,right=3mm,top=2mm,bottom=2mm}
\newtcolorbox{notebox}{colback=SoftGray,colframe=AccentTeal,boxrule=0.6pt,arc=2mm,left=3mm,right=3mm,top=2mm,bottom=2mm}
\newcommand{\DocTitle}[3]{%
\begin{center}
{\LARGE \bfseries \color{TitleBlue} #1}\\[0.5em]
{\large \color{AccentTeal} #2}\\[0.8em]
{\normalsize #3}
\end{center}
\vspace{0.5em}}
\newcommand{\ProjectTOC}{\vspace{0.5em}{\hypersetup{linkcolor=TitleBlue}\tableofcontents}\vspace{0.8em}}
\allowdisplaybreaks

\usepackage{tikz}
\usetikzlibrary{arrows.meta,positioning,calc,fit}
\usepackage{caption}
\usepackage{subcaption}
\usepackage{multirow}
\usepackage{makecell}
\usepackage{upquote}
\usepackage{array}

\definecolor{WarningRed}{HTML}{9B2C2C}
\definecolor{WarmGold}{HTML}{B7791F}
\definecolor{PaleGold}{HTML}{FFF8E6}
\definecolor{PaleRed}{HTML}{FFF1F1}
\definecolor{DeepGreen}{HTML}{276749}
\definecolor{PaleGreen}{HTML}{F0FFF4}
\definecolor{DeepPurple}{HTML}{5B3A8A}
\definecolor{PalePurple}{HTML}{F6F0FF}

\newtcolorbox{definitionbox}[1][]{
  colback=SoftTeal,colframe=AccentTeal,title={Definition#1},
  fonttitle=\bfseries,boxrule=0.7pt,arc=2mm,
  left=3mm,right=3mm,top=2mm,bottom=2mm}
\newtcolorbox{warningbox}[1][]{
  colback=PaleRed,colframe=WarningRed,title={Important distinction#1},
  fonttitle=\bfseries,boxrule=0.7pt,arc=2mm,
  left=3mm,right=3mm,top=2mm,bottom=2mm}
\newtcolorbox{checkpointbox}[1][]{
  colback=PaleGold,colframe=WarmGold,title={Learning checkpoint#1},
  fonttitle=\bfseries,boxrule=0.7pt,arc=2mm,
  left=3mm,right=3mm,top=2mm,bottom=2mm}
\newtcolorbox{conventionbox}[1][]{
  colback=SoftBlue,colframe=TitleBlue,title={Frozen project convention#1},
  fonttitle=\bfseries,boxrule=0.9pt,arc=2mm,
  left=3mm,right=3mm,top=2mm,bottom=2mm}
\newtcolorbox{practicebox}[1][]{
  colback=PaleGreen,colframe=DeepGreen,title={Recommended practice#1},
  fonttitle=\bfseries,boxrule=0.7pt,arc=2mm,
  left=3mm,right=3mm,top=2mm,bottom=2mm}
\newtcolorbox{advancedbox}[1][]{
  colback=PalePurple,colframe=DeepPurple,title={Advanced note#1},
  fonttitle=\bfseries,boxrule=0.7pt,arc=2mm,
  left=3mm,right=3mm,top=2mm,bottom=2mm}

\newcommand{\dd}{\mathrm{d}}
\newcommand{\pd}{\partial}
\newcommand{\vect}[1]{\bm{#1}}
\newcommand{\mat}[1]{\bm{#1}}
\newcommand{\average}[1]{\left\langle #1\right\rangle}
\newcommand{\bracket}[2]{\left[#1,#2\right]}
\newcommand{\Lap}{\nabla_\perp^2}
\newcommand{\Rpsi}{\mathcal{R}_{\psi}}
\newcommand{\Romega}{\mathcal{R}_{\omega}}
\newcommand{\Ldata}{\mathcal{L}_{\mathrm{data}}}
\newcommand{\Lpde}{\mathcal{L}_{\mathrm{PDE}}}
\newcommand{\Lic}{\mathcal{L}_{\mathrm{IC}}}
\newcommand{\Lbc}{\mathcal{L}_{\mathrm{BC}}}
\newcommand{\Lglobal}{\mathcal{L}_{\mathrm{global}}}
\newcommand{\Lgauge}{\mathcal{L}_{\mathrm{gauge}}}
\newcommand{\Ltotal}{\mathcal{L}_{\mathrm{total}}}
\newcommand{\MSE}{\operatorname{MSE}}
\newcommand{\RMSE}{\operatorname{RMSE}}
\newcommand{\RelLtwo}{\operatorname{RelL2}}
\newcommand{\FFT}{\mathcal{F}}
\newcommand{\IFFT}{\mathcal{F}^{-1}}
\newcommand{\Ptwo}{\mathcal{P}_{2/3}}
\newcommand{\code}[1]{\texttt{#1}}

\pagestyle{fancy}
\fancyhf{}
\lhead{PINN\_M3DC1}
\rhead{Module 5 Physics Notes}
\cfoot{\thepage}
\fancypagestyle{plain}{%
  \fancyhf{}%
  \lhead{PINN\_M3DC1}%
  \rhead{Module 5 Physics Notes}%
  \cfoot{\thepage}%
}

\begin{document}
\hypersetup{pageanchor=false}

\begin{titlepage}
\centering
\vspace*{1.0cm}
{\Large\color{AccentTeal} PINN\_M3DC1}\par
\vspace{0.65cm}
{\Huge\bfseries\color{TitleBlue} Module 5}\par
\vspace{0.25cm}
{\LARGE\bfseries\color{TitleBlue} Physics-Informed Full-Field\\Neural Operator}\par
\vspace{0.8cm}
{\large A textbook-style development of a reduced-MHD physics-informed neural operator on periodic numerical fields}\par
\vfill
\begin{tcolorbox}[width=0.92\textwidth,colback=SoftBlue,colframe=TitleBlue,arc=2mm]
\centering
\textbf{Learning task}\par
\vspace{0.35em}
Retain the Module~4 full-field operator map
\[
\boxed{
\left(\psi_0(x,y),U_0(x,y),\eta,\nu,t\right)
\longmapsto
\left(\widehat\psi(x,y,t),\widehat U(x,y,t)\right)
}
\]
and train the same operator family using the identical labeled field pairs together with the reduced-MHD equations, the initial-state identity, periodic-grid structure, the stream-function gauge, and optional global balance laws.
\end{tcolorbox}
\vfill
{\normalsize Version 1.1 \quad | \quad 11 August 2026}\par
\vspace{0.25cm}
{\small Physics and scientific-machine-learning note for the synthetic reduced-MHD proof of concept}\par
\end{titlepage}

\hypersetup{pageanchor=true}
\pagenumbering{roman}

\section*{Preface: from a data-only neural operator to a physics-informed neural operator}

Module~4 replaced the earlier pointwise coordinate network with a full-field Fourier neural operator (FNO). The data-only model receives complete numerical arrays for the initial magnetic-flux and stream-function fields, together with scalar transport coefficients and a requested target time. It returns complete numerical arrays for the target fields. Module~5 keeps that same scientific prediction task and the same primary FNO architecture, but now allows the reduced-MHD equations and other physical constraints to affect the learned weights.

The conceptual change is therefore
\[
\text{data-only neural operator}
\quad\longrightarrow\quad
\text{physics-informed neural operator}.
\]
In these notes, the abbreviation \textbf{PINO} means \emph{physics-informed neural operator}. The historical project name \code{PINN\_M3DC1} is retained, but the primary Module~5 learning architecture is no longer the coordinate-MLP PINN described in Version~1.0 of these notes.

This architectural revision changes how the physics residual is evaluated. The operator already produces an entire periodic spatial field on the common grid. Spatial derivatives can therefore be calculated by the same Fourier differentiation machinery used by the verified Module~2 pseudo-spectral solver. The scalar target time remains a differentiable network input, so the required time derivative is obtained by automatic differentiation with respect to that scalar. Thus the primary residual graph uses
\[
\boxed{\text{AD in time} \; + \; \text{spectral differentiation in space}.}
\]
The mathematical vorticity equation still contains fourth spatial derivatives of $U$, but the implementation no longer needs four nested automatic-differentiation operations with respect to coordinate inputs. Instead it applies the spectral Laplacian twice to the predicted full field.

\begin{warningbox}[\,: numerical fields are not images]
The PINO consumes and predicts numerical scalar arrays tied to a declared periodic grid. Contour plots, field-line renderings, quiver plots, color maps, PNG files, screenshots, and other rasterized figures are not learning inputs or targets. Physics residuals are evaluated from the numerical arrays themselves.
\end{warningbox}

\begin{warningbox}[\,: no automatic superiority]
Adding the correct equations does not guarantee a better surrogate. Physics losses can conflict with data losses, spectral truncation can hide current sheets, residual scaling can be poor, and the operator can satisfy the sampled equations while still predicting an inaccurate trajectory. Module~5 is a controlled experiment against Module~4, not an assumption that physics-informed training must win \cite{Wang2021Gradients,Krishnapriyan2021}.
\end{warningbox}

\begin{warningbox}[\,: claim boundary]
Module~5 remains a physics-informed surrogate of the project's independent two-dimensional, two-field reduced resistive-MHD model. It is not a validated surrogate of genuine M3D-C1 output. Real M3D-C1 arrays, geometry, closures, mesh mappings, and renewed verification are required for that stronger claim.
\end{warningbox}

\subsection*{How to read this note}

The main construction is
\[
\begin{gathered}
(\psi_0,U_0,\eta,\nu,t)
\longrightarrow \mathcal{G}_\theta
\longrightarrow (\widehat\psi_t,\widehat U_t)\\
\longrightarrow
\left\{\begin{array}{l}
\text{AD with respect to }t,\\
\text{Fourier derivatives in }x,y
\end{array}\right.
\longrightarrow
(\widehat J,\widehat\omega,\Rpsi,\Romega)\\
\longrightarrow
\text{physics-informed objective}
\longrightarrow
\text{updated operator weights}.
\end{gathered}
\]
Sections~\ref{sec:role}--\ref{sec:architecture} define the operator and physical problem. Sections~\ref{sec:residuals}--\ref{sec:derivatives} construct the residual graph. Sections~\ref{sec:populations}--\ref{sec:objective} define the training information and losses. Sections~\ref{sec:sampling}--\ref{sec:diagnostics} address sampling, optimization, and failure analysis. Sections~\ref{sec:experiments}--\ref{sec:contract} freeze the controlled experiment and implementation handoff.

\subsection*{Learning objectives}

After working through this module, the reader should be able to:
\begin{enumerate}[label=\arabic*.]
  \item state exactly what turns the Module~4 data-only FNO into a physics-informed neural operator;
  \item write the full-field map $(\psi_0,U_0,\eta,\nu,t)\mapsto(\widehat\psi_t,\widehat U_t)$;
  \item explain why $A_0$ remains case metadata but is not duplicated as a primary operator input;
  \item reproduce the project signs for $J$, $\omega$, $\Rpsi$, and $\Romega$;
  \item distinguish a coordinate PINN, a time-stepper, an FNO, and a PINO;
  \item explain why spatial derivatives are computed spectrally in the primary full-field formulation;
  \item explain why time derivatives are obtained by automatic differentiation with respect to scalar $t$;
  \item distinguish the mathematical derivative order of the PDE from the implementation depth of automatic differentiation;
  \item construct the solver-aligned dealiased Poisson brackets used in the physics residual;
  \item explain why inverse field scaling and correct time-scaling must occur inside the differentiable physics graph;
  \item distinguish labeled operator pairs from unlabeled physics contexts;
  \item explain why one full-grid physics context supplies many residual values but no new field labels;
  \item construct data, PDE, initial-condition, gauge, and optional global losses;
  \item explain why periodicity is structural in the primary periodic-grid FNO and therefore needs no seam penalty;
  \item design case-balanced physics-context sampling without leaking test trajectories;
  \item diagnose loss imbalance using per-term gradient magnitudes and gradient alignment;
  \item explain why low residual does not prove accurate reconnection physics;
  \item design warm-started and from-scratch comparisons with Module~4;
  \item state the current limitation imposed by the narrow family $U_0=A_0(\cos x-\cos y)$; and
  \item state precisely what Module~5 can and cannot establish about M3D-C1.
\end{enumerate}

\ProjectTOC
\clearpage
\pagenumbering{arabic}

\section{Scientific role of Module 5}
\label{sec:role}

\subsection{The controlled question}
Let $\mathcal{G}_{\rm data}$ denote the frozen Module~4 data-only FNO and $\mathcal{G}_{\rm phys}$ the Module~5 physics-informed operator. Both answer the same full-field question:
\begin{equation}
(\psi_0,U_0,\eta,\nu,t)\mapsto(\psi_t,U_t).
\end{equation}
The controlled scientific question is
\begin{quote}
At the same labeled full-field data budget and on the same untouched physical cases, does adding the declared reduced-MHD information improve field accuracy, derivative accuracy, current-sheet fidelity, reconnection timing, data efficiency, or physical consistency enough to justify the additional compute and optimization complexity?
\end{quote}
The comparison is meaningful only if the underlying prediction task, train/validation/test case split, labeled operator pairs, preprocessing, and baseline architecture remain controlled.

\subsection{What physics-informed means operationally}
\begin{definitionbox}
A model is \textbf{physics-informed during training} when a declared physical equation, condition, symmetry, conservation law, constitutive relation, or physical structural constraint influences the objective or representation that determines the trained parameters.
\end{definitionbox}

Module~5 can use five information channels:
\begin{enumerate}[label=\arabic*.]
  \item \textbf{labeled field data}: the exact Module~4 input-target operator pairs;
  \item \textbf{local dynamics}: reduced-MHD residuals evaluated on predicted full fields at unlabeled physics contexts;
  \item \textbf{initial-state identity}: at $t=0$, the output must equal the supplied initial fields;
  \item \textbf{periodicity and gauge}: periodic-grid structure and the convention $\average U=0$; and
  \item \textbf{global physics}: optional integrated balances such as energy decay and mean-flux conservation.
\end{enumerate}
Only the first channel contains reference target fields at $t>0$. The other channels constrain admissible predictions without revealing the hidden target arrays.

\subsection{From PINN to PINO}
A classical coordinate PINN represents $q(x,y,t)$ directly as a scalar-valued neural function and differentiates that function with respect to coordinates. The revised Module~5 represents a function-to-function map. The network consumes $q_0(x,y)$ and returns an entire $q_t(x,y)$ array.

\begin{longtable}{p{0.20\textwidth}p{0.34\textwidth}p{0.36\textwidth}}
\toprule
Formulation & Typical map & Primary derivative strategy \\
\midrule
\endfirsthead
\toprule
Formulation & Typical map & Primary derivative strategy \\
\midrule
\endhead
Coordinate PINN & $(x,y,t,\mu)\mapsto q$ & AD with respect to spatial coordinates and time \\
Autoregressive neural time-stepper & $q_n\mapsto q_{n+1}$ & Usually discrete-step consistency or rollout losses \\
Data-only FNO & $(q_0,\mu,t)\mapsto q_t$ & No PDE derivative required during training \\
Module~5 PINO & $(q_0,\mu,t)\mapsto q_t$ & AD in scalar time; spectral derivatives on output fields \\
\bottomrule
\end{longtable}

The word \emph{operator} does not mean the model is exact on arbitrary functions. It means the learning object is a map between fields/functions. The actual generalization domain is still determined by the training distribution.

\subsection{Hybrid, data-free, and data-only limits}
The composite objective permits three regimes:
\begin{align}
\text{data-only: }&w_{\rm data}>0,\quad w_{\rm PDE}=w_{\rm IC}=w_{\rm gauge}=w_{\rm global}=0,\\
\text{hybrid PINO: }&w_{\rm data}>0,\quad \text{at least one physics term active},\\
\text{data-free physics operator: }&w_{\rm data}=0,\quad \text{PDE and trajectory-defining conditions active}.
\end{align}
The primary project comparison uses the hybrid regime. A data-free operator is scientifically interesting but answers a different question about information and computational cost.

\subsection{A PINO is not the reference solver}
The Module~2 solver advances $(\psi,\omega)$ through explicit RK4 stages and reconstructs $U$ and $J$ using Fourier operators. The Module~5 operator predicts a requested time directly from the initial fields. It does not need to march through earlier saved states.

Even when the physics loss uses the same Fourier derivative convention as the solver, the neural operator is not simply a rewritten RK4 integrator. Its state representation, optimization error, spectral mode bottlenecks, and time-conditioning mechanism are different.

\section{Frozen physical problem and operator contract}
\label{sec:physics}

\subsection{Domain and fields}
The spatial domain is
\begin{equation}
\Omega=[0,L_x)\times[0,L_y),\qquad L_x=L_y=2\pi
\end{equation}
in the version-1 baseline, with $t\in[0,T]$. The primary physical fields are magnetic flux $\psi$ and flow stream function $U$. Derived fields are
\begin{equation}
J=-\Lap\psi,
\qquad
\omega=\Lap U.
\label{eq:derived}
\end{equation}
The in-plane magnetic and velocity fields are
\begin{align}
\vect B_\perp&=(\psi_y,-\psi_x),\\
\vect v_\perp&=(-U_y,U_x).
\end{align}

\subsection{Poisson bracket}
For scalar fields $f$ and $g$,
\begin{equation}
\bracket{f}{g}=f_xg_y-f_yg_x.
\label{eq:bracket}
\end{equation}
Under the project convention, $[U,f]=\vect v_\perp\cdot\nabla f$.

\subsection{Evolution equations}
The frozen reduced-MHD model is
\begin{align}
\psi_t+\bracket{U}{\psi}&=\eta\Lap\psi,
\label{eq:flux}\qquad\\
\omega_t+\bracket{U}{\omega}&=\bracket{J}{\psi}+\nu\Lap\omega,
\label{eq:vort}\qquad\\
J&=-\Lap\psi,\qquad \omega=\Lap U.
\end{align}
All variables are nondimensional under Module~1 normalization.

\subsection{Initial condition and $A_0$}
Version~1 uses
\begin{equation}
\psi_0(x,y)=\sin x\sin y,
\qquad
U_0(x,y)=A_0(\cos x-\cos y).
\end{equation}
The operator receives the complete $\psi_0$ and $U_0$ arrays. Therefore $A_0$ is already encoded in $U_0$ and is not duplicated as a primary input channel. It remains mandatory metadata because it identifies the solver case and supports consistency checks.

\begin{warningbox}[\,: narrow initial-state family]
Although the network receives a complete initial field, version~1 still varies only the amplitude of one fixed stream-function shape. Holding out one $A_0$ tests amplitude generalization within this family, not generalization to arbitrary magnetic topology or arbitrary flow morphology.
\end{warningbox}

\subsection{Exact operator map}
For case $c$ and requested time $t$,
\begin{equation}
\boxed{
\mathcal{G}_\theta:
\left(\psi_0^{(c)},U_0^{(c)},\eta_c,\nu_c,t\right)
\longmapsto
\left(\widehat\psi^{(c)}_t,\widehat U^{(c)}_t\right).
}
\label{eq:operator}
\end{equation}
Each field is a complete $N_x\times N_y$ numerical array on the common periodic grid.

\subsection{Initial, periodic, and gauge information}
The time-zero operator must satisfy
\begin{equation}
\mathcal{G}(\psi_0,U_0,\eta,\nu,0)=(\psi_0,U_0).
\label{eq:icidentity}
\end{equation}
All fields are interpreted periodically in $x$ and $y$. The stream-function gauge is
\begin{equation}
\average U(t)=\frac{1}{|\Omega|}\int_\Omega U\,\dd A=0.
\end{equation}
The spatial mean of $\psi$ is fixed by the initial condition and preserved by the periodic PDE.

\section{The matched Module 5 neural operator architecture}
\label{sec:architecture}

\subsection{Architecture inherited from Module 4}
The primary comparison keeps the Module~4 two-dimensional FNO family. The input tensor contains five channels:
\begin{equation}
\mat X(x_i,y_j)=
\begin{bmatrix}
\widetilde\psi_0(x_i,y_j)\\
\widetilde U_0(x_i,y_j)\\
\widetilde\eta\\
\widetilde\nu\\
\widetilde t
\end{bmatrix}.
\end{equation}
The three scalar context values are broadcast as spatially constant channels. A lifting layer maps these channels into a wider latent representation. Repeated FNO blocks combine learned Fourier-mode mixing with local pointwise channel mixing, followed by a smooth nonlinearity. A projection head returns two output channels.

\subsection{One FNO block}
Let $v^{(\ell)}(x,y)\in\mathbb{R}^{d}$ be the latent field at layer $\ell$. A schematic block is
\begin{equation}
v^{(\ell+1)}
=\sigma\!\left(
W^{(\ell)}v^{(\ell)}
+\IFFT\left(R^{(\ell)}(k_x,k_y)\,\FFT[v^{(\ell)}]\right)
+b^{(\ell)}
\right),
\label{eq:fno_block}
\end{equation}
where $R^{(\ell)}$ contains learned complex channel-mixing weights on retained Fourier modes and is zero outside the learned spectral band.

\subsection{Why the architecture stays matched}
If Module~5 were given a larger network, more Fourier modes, a different activation, or a different preprocessing pipeline, an observed improvement could not be attributed cleanly to physics information. The primary comparison therefore uses the same FNO width, layer count, mode count, activation, projection, padding rule, and tensor contract selected for Module~4.

\subsection{Periodicity is structural}
The physical grid is half-open and periodic, and the FNO spectral branch uses discrete Fourier transforms on that grid. There is no duplicated $x=L_x$ or $y=L_y$ endpoint. The predicted array is interpreted as one period of a periodic discrete field.

For the primary architecture,
\begin{equation}
\Lbc=0.
\end{equation}
A separate boundary-seam loss is unnecessary unless a deliberately nonperiodic architecture is tested as an ablation.

\subsection{Smoothness requirement changes}
The old coordinate PINN required an activation differentiable through fourth order because spatial derivatives were obtained through nested AD. The full-field PINO computes spatial derivatives spectrally, so the activation does not need four classical coordinate derivatives for that purpose.

However, the output must remain differentiable with respect to the scalar time input. A smooth activation such as GELU is therefore retained for the matched Module~4/5 baseline. Non-smooth alternatives would require explicit validation of the time derivative graph.

\subsection{Why the FNO mode count still matters to physics}
Current and vorticity amplify high wave numbers:
\begin{equation}
\widehat J_{\vect k}=k^2\widehat\psi_{\vect k},
\qquad
\widehat\omega_{\vect k}=-k^2\widehat U_{\vect k}.
\end{equation}
An operator can achieve low field MSE while smoothing a thin current sheet. The physics residual does not remove this architectural bottleneck automatically. Mode-count and derivative-sensitive validation remain essential.

\section{From full-field predictions to reduced-MHD residuals}
\label{sec:residuals}

\subsection{Residual definitions}
Move each evolution equation to the left:
\begin{align}
\Rpsi
&=\widehat\psi_t+\bracket{\widehat U}{\widehat\psi}
-\eta\Lap\widehat\psi,
\label{eq:Rpsi}\qquad\\
\Romega
&=\widehat\omega_t+\bracket{\widehat U}{\widehat\omega}
-\bracket{\widehat J}{\widehat\psi}
-\nu\Lap\widehat\omega,
\label{eq:Romega}
\end{align}
with
\begin{equation}
\widehat J=-\Lap\widehat\psi,
\qquad
\widehat\omega=\Lap\widehat U.
\end{equation}
An exact sufficiently smooth solution has zero residual everywhere.

\subsection{Fully expanded mathematical form}
The flux residual is
\begin{equation}
\Rpsi=\widehat\psi_t
+\widehat U_x\widehat\psi_y
-\widehat U_y\widehat\psi_x
-\eta(\widehat\psi_{xx}+\widehat\psi_{yy}).
\end{equation}
The vorticity residual is
\begin{align}
\Romega={}&\widehat\omega_t
+\widehat U_x\widehat\omega_y
-\widehat U_y\widehat\omega_x
-\widehat J_x\widehat\psi_y
+\widehat J_y\widehat\psi_x\\
&-\nu(\widehat\omega_{xx}+\widehat\omega_{yy}).
\end{align}
Mathematically, the final viscous term contains fourth spatial derivatives of $U$.

\subsection{Derivative-order audit}
\begin{longtable}{p{0.25\textwidth}p{0.28\textwidth}p{0.37\textwidth}}
\toprule
Residual quantity & Highest mathematical derivative & Primary implementation \\
\midrule
\endfirsthead
\toprule
Residual quantity & Highest mathematical derivative & Primary implementation \\
\midrule
\endhead
$\widehat\psi_t$ & first in time & AD with respect to scalar $t$ \\
$[\widehat U,\widehat\psi]$ & first in space & Fourier first derivatives + physical products \\
$\Lap\widehat\psi$ & second in space & Fourier multiplier $-k^2$ \\
$\widehat\omega$ & second in space & Fourier Laplacian of $\widehat U$ \\
$\widehat\omega_t$ & second in space + first in time & $\Lap\widehat U_t$ \\
$[\widehat U,\widehat\omega]$ & third in $U$ & spectral derivatives of $\widehat\omega$ \\
$[\widehat J,\widehat\psi]$ & third in $\psi$ & spectral derivatives of $\widehat J$ and $\widehat\psi$ \\
$\Lap\widehat\omega$ & fourth in $U$ & second spectral Laplacian \\
\bottomrule
\end{longtable}

\begin{checkpointbox}
The PDE has not become lower order. The vorticity equation still contains a biharmonic contribution in $U$. What changed is the computational route: Fourier multipliers evaluate those spatial derivatives directly on the predicted field instead of constructing four nested spatial AD graphs.
\end{checkpointbox}

\subsection{Residual of the continuous equation versus residual on the discrete grid}
A field represented on an $N_x\times N_y$ periodic grid is a finite numerical object. Spatial differentiation and nonlinear multiplication therefore require a discrete convention. The primary Module~5 residual uses the same Fourier wave numbers and rectangular two-thirds dealiasing convention verified in Module~2.

This is best described as a \emph{solver-aligned spectral strong residual}: it approximates the continuous reduced-MHD equation using the accepted periodic spectral representation while controlling quadratic aliasing. It is not the algebraic RK4 stage residual of the reference integrator.

\subsection{Why a small residual is not sufficient}
A low sampled residual does not by itself prove that the predicted trajectory is correct. The wrong initial state, wrong gauge, inadequate spectral resolution, residual normalization, or a narrow physics-context distribution can all produce misleadingly small averages. The final test therefore evaluates complete held-out fields and physical diagnostics in addition to residuals.

\section{Spectral spatial differentiation and dealiased nonlinear physics}
\label{sec:spectral}

\subsection{Fourier derivative operators}
For a periodic numerical field $f$, let $\widehat f_{mn}=\FFT[f]_{mn}$. With $k^2=k_x^2+k_y^2$,
\begin{align}
\FFT[f_x]&=ik_x\widehat f,\\
\FFT[f_y]&=ik_y\widehat f,\\
\FFT[\Lap f]&=-k^2\widehat f.
\end{align}
Therefore
\begin{align}
\FFT[J]&=k^2\widehat\psi,\\
\FFT[\omega]&=-k^2\widehat U,\\
\FFT[\Lap\omega]&=k^4\widehat U.
\end{align}
The zero mode is treated consistently with the Module~2 gauge and mean conventions.

\subsection{Dealiased Poisson bracket}
For resolved fields $f$ and $g$, define the solver-aligned bracket operator
\begin{equation}
\mathcal{B}_{2/3}(f,g)
=\IFFT\!\left[
\Ptwo\FFT\!\left(f_xg_y-f_yg_x\right)
\right].
\label{eq:dealiased_bracket}
\end{equation}
The rectangular projection $\Ptwo$ retains the same mode set used by Module~2.

The primary physics residual uses
\begin{equation}
\Rpsi=\widehat\psi_t+\mathcal{B}_{2/3}(\widehat U,\widehat\psi)-\eta\Lap\widehat\psi,
\end{equation}
\begin{equation}
\Romega=\widehat\omega_t+\mathcal{B}_{2/3}(\widehat U,\widehat\omega)
-\mathcal{B}_{2/3}(\widehat J,\widehat\psi)-\nu\Lap\widehat\omega.
\end{equation}

\subsection{Why dealias inside the residual}
If nonlinear products generate wave numbers beyond the grid's representable band, unresolved contributions can fold into lower modes and create false residual structure. Dealiasing only the final scalar loss is too late because the aliased product has already contaminated the local equation defect.

\subsection{Residual projection is a numerical contract}
The two-thirds projection changes the finite-grid residual norm being minimized. It does not change the continuous PDE written in Module~1. For reproducibility, every run must record the derivative convention, FFT normalization, retained mask, and whether the residual is projected before squaring.

\subsection{Alternative weak or multiresolution physics losses}
A weak-form residual, finite-volume residual, or coarse-grid residual can reduce sensitivity to unresolved high-frequency derivatives. Those are scientifically legitimate alternatives but define different physics losses. They belong in controlled ablations after the primary spectral strong-form implementation is verified.

\section{Time differentiation and the differentiable operator graph}
\label{sec:derivatives}

\subsection{Why time is different from space}
The FNO produces a field on a fixed spatial grid, so spectral spatial derivatives are natural. Time, however, is supplied as one scalar conditioning variable. The operator is intended to approximate a continuous family $\mathcal{G}_\theta(\cdot,t)$. Module~5 therefore differentiates the predicted fields with respect to that scalar.

\subsection{Automatic differentiation with respect to time}
For one grid point, the network output is a differentiable function of $t$ through the broadcast time channel. Automatic differentiation applies the chain rule through time preprocessing, lifting, FNO layers, and projection to obtain
\begin{equation}
\widehat\psi_t(x_i,y_j),\qquad \widehat U_t(x_i,y_j)
\end{equation}
for every grid point in the predicted field.

The vorticity time derivative can then be formed as
\begin{equation}
\widehat\omega_t=\Lap\widehat U_t,
\end{equation}
which commutes with time differentiation for the smooth represented field.

\subsection{AD is not finite differencing in time}
A centered finite difference would use
\begin{equation}
\widehat\psi_t(t)\approx
\frac{\widehat\psi(t+h)-\widehat\psi(t-h)}{2h}.
\end{equation}
That requires extra operator evaluations and introduces a time-step parameter $h$. The primary formulation instead differentiates the implemented continuous time-conditioned operator directly. Finite-difference time residuals may be tested later as a computational ablation.

\subsection{Scaling chain rule}
Suppose
\begin{equation}
\widetilde t=\frac{t-m_t}{s_t},
\qquad
\widehat\psi=m_\psi+s_\psi\widehat{\widetilde\psi}.
\end{equation}
Then
\begin{equation}
\widehat\psi_t
=\frac{s_\psi}{s_t}
\frac{\pd\widehat{\widetilde\psi}}{\pd\widetilde t}.
\label{eq:time_chain}
\end{equation}
An omitted factor changes the physical coefficient balance in the PDE.

\begin{practicebox}
The safest implementation exposes one differentiable function of the \emph{raw nondimensional} inputs. That wrapper performs stored preprocessing, evaluates the FNO, inverse-transforms the outputs, and returns physical nondimensional $\widehat\psi$ and $\widehat U$. The physics residual differentiates this wrapper with respect to raw $t$.
\end{practicebox}

\subsection{Transport-coefficient preprocessing}
The FNO may receive transformed coefficients such as $\log_{10}\eta$ and $\log_{10}\nu$, but the PDE residual must multiply diffusion terms by the physical nondimensional $\eta$ and $\nu$. The transformed channels condition the network; they do not replace the physical coefficients in Equations~\eqref{eq:Rpsi}--\eqref{eq:Romega}.

\subsection{Precision, memory, and time derivative cost}
A full-field time derivative requires differentiating many outputs with respect to one scalar input. Forward-mode AD can be attractive because the number of differentiated inputs is small, whereas reverse-mode backpropagation remains natural for the scalar training loss with respect to many network weights. Framework-specific choices must be benchmarked rather than assumed.

High-order spatial differentiation also amplifies numerical high-frequency content. Single- and double-precision residual comparisons should be performed on manufactured fields before production training.

\section{Training information populations}
\label{sec:populations}

\subsection{Labeled operator pairs}
The data population is identical to Module~4. For case $c$ and saved target time $t_n$,
\begin{equation}
X^d_{c,n}=(\psi_0^{(c)},U_0^{(c)},\eta_c,\nu_c,t_n),
\end{equation}
\begin{equation}
Y^d_{c,n}=(\psi_n^{(c)},U_n^{(c)}).
\end{equation}
These pairs contain reference target fields and determine $\Ldata$.

\subsection{Unlabeled physics contexts}
A physics context contains
\begin{equation}
X^f=(\psi_0,U_0,\eta,\nu,t_f)
\end{equation}
without a target field at $t_f$. The operator predicts a complete field, and Equations~\eqref{eq:Rpsi}--\eqref{eq:Romega} generate the training signal.

A physics time may lie between saved solver snapshots. This is one concrete way the equation provides information where no label is stored.

\subsection{One context gives a residual field, not thousands of labels}
For a $128\times128$ grid, one physics context produces $16{,}384$ values of $\Rpsi$ and $16{,}384$ values of $\Romega$. These are spatial equation defects, not independent target observations. The project separately reports:
\begin{itemize}
  \item number of independent training trajectories;
  \item number of labeled operator pairs;
  \item number of physics contexts evaluated;
  \item number of residual grid nodes evaluated; and
  \item wall-clock and accelerator cost.
\end{itemize}

\subsection{Initial-condition contexts}
At $t=0$, the supplied input fields are exact targets by definition. Initial contexts therefore use
\begin{equation}
(\psi_0,U_0,\eta,\nu,0)
\end{equation}
and compare the prediction directly with $(\psi_0,U_0)$.

\subsection{Gauge and global groups}
Because every prediction is already a complete spatial field, a gauge or global balance can be evaluated using the same grid and quadrature weights. No separate collection of arbitrary coordinate points is required.

\subsection{Physics contexts and the train/test boundary}
\begin{warningbox}[\,: unlabeled test leakage is still leakage]
If the scientific test is generalization to unseen physical cases, Module~5 may not use the test case's initial fields, exact $(\eta,\nu)$ combination, or target-time contexts in a physics loss during training. The absence of target labels does not make transductive use of the held-out case invisible.
\end{warningbox}

The primary experiment restricts physics contexts to training trajectories and their declared coefficient values. Physics training at new unsolved parameter combinations is a separate experiment and changes the meaning of ``unseen case.''

\section{Supervised full-field data loss}
\label{sec:data_loss}

\subsection{Identical Module 4 observations}
Module~5 inherits the immutable labeled pair manifest from Module~4. No additional solver target fields may be added to the physics-informed model for the primary comparison.

\subsection{Channel-balanced field error}
For pair $(c,n)$, standardized complete-field errors are
\begin{align}
e_{\psi,c,n,i,j}&=\widehat{\widetilde\psi}_{c,n,i,j}-\widetilde\psi_{c,n,i,j},\\
e_{U,c,n,i,j}&=\widehat{\widetilde U}_{c,n,i,j}-\widetilde U_{c,n,i,j}.
\end{align}
Define
\begin{align}
E_{\psi,c,n}&=\frac{1}{N_xN_y}\sum_{i,j}e_{\psi,c,n,i,j}^2,\\
E_{U,c,n}&=\frac{1}{N_xN_y}\sum_{i,j}e_{U,c,n,i,j}^2.
\end{align}
After separate output standardization, the pair loss is
\begin{equation}
\ell^d_{c,n}=E_{\psi,c,n}+E_{U,c,n}.
\end{equation}

\subsection{Case- and time-balanced risk}
If $\mathcal{B}_C$ is a set of sampled training cases and $\mathcal{B}_t(c)$ the selected labeled times within case $c$,
\begin{equation}
\boxed{
\Ldata=
\frac{1}{|\mathcal{B}_C|}
\sum_{c\in\mathcal{B}_C}
\frac{1}{|\mathcal{B}_t(c)|}
\sum_{n\in\mathcal{B}_t(c)}\ell^d_{c,n}.
}
\label{eq:data_loss}
\end{equation}
This is exactly the primary Module~4 data loss.

\begin{conventionbox}[\,: fairness anchor]
The Module~5 primary comparison uses the same labeled pairs, output scaling, channel weights, case balancing, and target-time balancing as Module~4. Physics information is added; labeled information is not.
\end{conventionbox}

\section{PDE residual loss}
\label{sec:pde_loss}

\subsection{Residual normalization}
The two equations can have different numerical scales. Let $S_{R_\psi}>0$ and $S_{R_\omega}>0$ be residual scales estimated from training information only. Define
\begin{equation}
r_\psi=\frac{\Rpsi}{S_{R_\psi}},
\qquad
r_\omega=\frac{\Romega}{S_{R_\omega}}.
\end{equation}
The entire physical residual is formed before this normalization.

\begin{warningbox}[\,: never normalize PDE terms independently]
Do not separately standardize $\psi_t$, advection, Lorentz forcing, and diffusion before adding them. The physical equation depends on cancellation among terms with fixed relative coefficients. Independent normalization changes the equation.
\end{warningbox}

\subsection{Full-grid residual norm}
For one physics context $f$,
\begin{align}
E_{R_\psi,f}&=\frac{1}{N_xN_y}\sum_{i,j}r_{\psi,f,i,j}^2,\\
E_{R_\omega,f}&=\frac{1}{N_xN_y}\sum_{i,j}r_{\omega,f,i,j}^2.
\end{align}
The primary PDE loss is
\begin{equation}
\Lpde=
\frac{1}{|\mathcal{B}_f|}
\sum_{f\in\mathcal{B}_f}
\left(E_{R_\psi,f}+E_{R_\omega,f}\right),
\label{eq:pde_loss}
\end{equation}
with case-balanced sampling of contexts.

\subsection{Area averages can hide current sheets}
A current sheet may occupy a small fraction of the domain while controlling reconnection. Uniform MSE can therefore be small even when peak residual near the sheet is large. Training logs should record at least
\begin{equation}
\|r_\psi\|_\infty,\qquad \|r_\omega\|_\infty
\end{equation}
in addition to mean-square values.

A structure-weighted residual is a useful ablation, but the primary physics loss begins with uniform spatial weighting so its norm is unambiguous.

\subsection{Reference residual floor}
Accepted Module~2 data do not satisfy the continuous PDE with exact zero numerical defect. Truncation, time integration, and saved representation create a finite residual floor under any chosen diagnostic operator. Before judging a PINO residual, the project evaluates the same spectral residual convention on accepted reference snapshots and records the resulting scale.

\section{Initial-condition enforcement}
\label{sec:ic}

\subsection{The operator identity at zero time}
Because the complete initial field is already supplied to the model, the correct time-zero condition is especially simple:
\begin{equation}
\widehat\psi(x,y,0)=\psi_0(x,y),
\qquad
\widehat U(x,y,0)=U_0(x,y).
\end{equation}

\subsection{Soft initial-condition loss}
For a batch of initial contexts,
\begin{equation}
\Lic=
\average{
\frac{\|\widehat\psi_0-\psi_0\|_2^2}{S_\psi^2|\Omega|}
+
\frac{\|\widehat U_0-U_0\|_2^2}{S_U^2|\Omega|}
}_{\rm IC\ contexts}.
\label{eq:ic_loss}
\end{equation}
This is treated as physics/trajectory information rather than an additional interior solver snapshot.

\subsection{Hard initial-condition operator ansatz}
A future ablation can impose
\begin{align}
\widehat\psi_t&=\psi_0+g(t)\,\mathcal{N}_{\psi,\theta}(\psi_0,U_0,\eta,\nu,t),\\
\widehat U_t&=U_0+g(t)\,\mathcal{N}_{U,\theta}(\psi_0,U_0,\eta,\nu,t),
\end{align}
with $g(0)=0$. Then the initial condition holds for every parameter value and every network weight.

The primary Module~5 model uses the soft loss first because it keeps the Module~4 architecture exactly matched. Hard enforcement changes the output parameterization and therefore belongs in a declared ablation.

\subsection{Initial derivative diagnostics}
A field can match $\psi_0$ and $U_0$ approximately while producing inaccurate $J_0$ or $\omega_0$. The project therefore diagnoses
\begin{equation}
\widehat J_0=-\Lap\widehat\psi_0,
\qquad
\widehat\omega_0=\Lap\widehat U_0,
\end{equation}
but does not add target-current or target-vorticity losses to the primary objective.

\section{Periodicity, gauge, and mean constraints}
\label{sec:gauge}

\subsection{Periodicity through the full-field representation}
The network output is a field on the half-open periodic grid and the primary spatial physics operators are Fourier operators. The representation is therefore periodically extended by construction. No separate matched boundary-point loader is required.

\subsection{Why the gauge remains necessary}
Adding a spatial constant $C(t)$ to $U$ changes neither velocity nor vorticity:
\begin{equation}
U'(x,y,t)=U(x,y,t)+C(t).
\end{equation}
The local reduced-MHD residual is blind to this spatial constant. The project fixes
\begin{equation}
\average U(t)=0.
\end{equation}

\subsection{Soft gauge loss}
On the full grid,
\begin{equation}
\Lgauge=
\average{
\left(
\frac{\frac{1}{N_xN_y}\sum_{i,j}\widehat U_{ij}}{S_U}
\right)^2
}_{\rm gauge\ contexts}.
\label{eq:gauge_loss}
\end{equation}
A differentiable exact mean-subtraction projection is possible, but it changes the output wrapper and should be tested separately.

\subsection{Mean magnetic flux}
Integrating the flux equation over the periodic domain gives
\begin{equation}
\frac{\dd}{\dd t}\average\psi=0.
\end{equation}
An optional mean-flux loss is
\begin{equation}
\mathcal{L}_{\langle\psi\rangle}
=\average{
\left(
\frac{\average{\widehat\psi(t)}-\average{\psi_0}}{S_\psi}
\right)^2
}.
\end{equation}
It is useful but redundant with perfect local physics and conditions; version~1 treats it as an optional global term.

\section{Global balance constraints}
\label{sec:global}

\subsection{Magnetic, kinetic, and total energy}
Define
\begin{align}
E_B(t)&=\frac12\int_\Omega |\nabla\psi|^2\,\dd A,\\
E_K(t)&=\frac12\int_\Omega |\nabla U|^2\,\dd A,\\
E(t)&=E_B(t)+E_K(t).
\end{align}
For smooth periodic solutions,
\begin{equation}
\boxed{
\frac{\dd E}{\dd t}
=-\eta\int_\Omega J^2\,\dd A
-\nu\int_\Omega \omega^2\,\dd A.
}
\label{eq:energy_balance}
\end{equation}

\subsection{Sketch of the magnetic-energy derivation}
Using $J=-\Lap\psi$ and periodic integration by parts,
\begin{equation}
\frac{\dd E_B}{\dd t}
=\int_\Omega J\,\psi_t\,\dd A.
\end{equation}
Insert the flux equation:
\begin{equation}
\frac{\dd E_B}{\dd t}
=-\int_\Omega J[U,\psi] \,\dd A
-\eta\int_\Omega J^2\,\dd A.
\end{equation}
The first term exchanges energy with the flow.

\subsection{Sketch of the kinetic-energy derivation}
Using $\omega=\Lap U$,
\begin{equation}
\frac{\dd E_K}{\dd t}
=-\int_\Omega U\,\omega_t\,\dd A.
\end{equation}
Substitution of the vorticity equation and the periodic bracket identities gives the opposite magnetic-flow exchange term plus
\begin{equation}
-\nu\int_\Omega\omega^2\,\dd A.
\end{equation}
Adding the two balances produces Equation~\eqref{eq:energy_balance}.

\subsection{Instantaneous energy-balance residual}
Because the operator is differentiable with respect to $t$, define
\begin{equation}
R_E(t)=
\frac{\dd \widehat E}{\dd t}
+\eta\int_\Omega \widehat J^2\,\dd A
+\nu\int_\Omega \widehat\omega^2\,\dd A.
\end{equation}
The derivative $\dd\widehat E/\dd t$ can be obtained by AD through the global energy calculation.

\subsection{Mean-square flux balance}
Let
\begin{equation}
A(t)=\frac12\int_\Omega\psi^2\,\dd A.
\end{equation}
Then
\begin{equation}
\frac{\dd A}{\dd t}=-2\eta E_B.
\label{eq:mean_square_flux}
\end{equation}
This supplies another optional scalar global defect.

\subsection{Why global losses are optional}
A global integral is much less informative than a local PDE. Many incorrect fields share the same total energy. Global balances are therefore secondary regularizers and diagnostics, not replacements for $\Rpsi$ and $\Romega$.

\section{Composite physics-informed objective}
\label{sec:objective}

\subsection{Top-level loss}
The general objective is
\begin{equation}
\boxed{
\Ltotal
=w_d\Ldata
+w_f\Lpde
+w_0\Lic
+w_g\Lgauge
+w_G\Lglobal.
}
\label{eq:total_loss}
\end{equation}
For the primary periodic FNO, $\Lbc=0$ structurally.

\subsection{The primary information ladder}
To identify which physics information helps, use the ladder
\begin{align*}
\mathcal{M}_0 &: \Ldata,\\
\mathcal{M}_1 &: \Ldata+\Lic,\\
\mathcal{M}_2 &: \Ldata+\Lic+\Lpde^{(\psi)},\\
\mathcal{M}_3 &: \Ldata+\Lic+\Lpde^{(\psi)}+\Lpde^{(\omega)}+\Lgauge,\\
\mathcal{M}_4 &: \mathcal{M}_3+\Lglobal.
\end{align*}
Module~4 is $\mathcal{M}_0$. The primary Module~5 endpoint is $\mathcal{M}_3$; global terms are a controlled extension rather than a hidden ingredient.

\subsection{A weighted sum is a multi-objective compromise}
Different losses can prefer different parameter updates. A small scalar loss does not imply a small influence on the weights, and equal loss values do not imply balanced gradients.

\subsection{Residual-first scaling hierarchy}
A robust order is:
\begin{enumerate}
  \item construct every physical residual in raw nondimensional variables;
  \item normalize whole residual equations by declared characteristic scales;
  \item make each component loss an interpretable average;
  \item inspect gradient norms with unit top-level weights; and
  \item choose or adapt the top-level weights from training/validation information only.
\end{enumerate}

\subsection{Fixed versus adaptive weights}
Fixed weights are easiest to reproduce. Adaptive strategies can equalize gradient magnitudes, enforce constraints with multipliers, or respond to training progress. They may improve optimization but add another algorithm whose behavior must be reported.

\subsection{Weight-selection leakage}
Loss weights, residual scales, and stopping rules must not be tuned on the final test cases. Physics-informed training is especially vulnerable to this mistake because many apparently reasonable weights can be tried.

\section{Physics-context sampling}
\label{sec:sampling}

\subsection{Sampling now acts mainly over cases and time}
The primary spatial residual is evaluated on the full periodic grid because FFT derivatives require the complete field. The principal sampling choices are therefore:
\begin{itemize}
  \item which training cases enter a batch;
  \item which labeled target times enter the data loss;
  \item which unlabeled times enter the physics loss; and
  \item how those times are distributed across early, current-sheet, reconnection, and post-event phases.
\end{itemize}

\subsection{Continuous unlabeled time sampling}
Within a training case, $t_f$ can be sampled continuously over the declared interval rather than restricted to saved output times. This lets the equation constrain the interpolating operator between reference snapshots.

\subsection{Do not silently expand the parameter domain}
Because the governing PDE is known for arbitrary positive $\eta$ and $\nu$, it is tempting to sample new coefficients between solver cases. That is a legitimate physics-only extension, but it changes the training support. If the final test contains those same coefficients, the experiment is no longer strict unseen-parameter generalization.

The version-1 primary experiment therefore uses only training-case coefficient values for physics contexts.

\subsection{Time curriculum}
Current-sheet formation creates progressively smaller spatial scales. A time curriculum can begin with early smooth dynamics and extend the maximum physics time after the model stabilizes. The curriculum boundary and schedule must be logged.

\subsection{Residual-adaptive time sampling}
A candidate pool of times can be evaluated periodically, with more physics contexts allocated where dense validation residuals are large. This is analogous to residual-adaptive refinement, but the adaptation is primarily in case-time space rather than arbitrary coordinate space.

\subsection{Structure-aware spatial weighting}
Although FFT derivatives require a full field, the residual MSE may be reweighted after the residual map is computed. Regions of high predicted or reference current can receive additional weight. Using reference current introduces extra target-derived information and must therefore be matched carefully with Module~4 or labeled as an enhanced physics/data ablation.

\section{Optimization and training curriculum}
\label{sec:training}

\subsection{Why warm-starting is natural}
Module~4 has already trained the exact same operator architecture on the exact same labeled data. A clean way to ask what the physics does is therefore to initialize Module~5 from the frozen Module~4 checkpoint and continue optimization with physics terms active.

This warm start does not add labels, but it adds training compute. Any compute-efficiency claim must include the Module~4 pretraining cost.

\subsection{From-scratch training remains necessary as an ablation}
A physics-informed model trained from random initialization answers whether the composite objective can train the operator directly. Comparing warm-started and from-scratch runs separates representational benefit from optimization assistance.

\subsection{Stage 0: operator and residual unit tests}
Before learned training:
\begin{itemize}
  \item verify FFT wave numbers and Laplacians on analytic Fourier modes;
  \item verify $J=-\Lap\psi$ and $\omega=\Lap U$ signs;
  \item verify bracket antisymmetry and $[f,f]=0$;
  \item verify two-thirds projection matches Module~2;
  \item verify raw-time AD against analytic manufactured time dependence;
  \item verify scaling invariance of the reconstructed physical residual; and
  \item verify all active losses backpropagate finite gradients to FNO weights.
\end{itemize}

\subsection{Stage 1: matched Module 4 warm start}
Load the selected Module~4 checkpoint and confirm identical predictions and $\Ldata$ before any physics update.

\subsection{Stage 2: initial identity and flux equation}
Activate $\Lic$ and $\Rpsi$ first. The flux equation reaches only second spatial derivatives and is the simpler residual branch.

\subsection{Stage 3: full vorticity physics}
Activate $\Romega$, derived $J$ and $\omega$, and the gauge loss. Monitor high-wave-number spectra and nonfinite values closely.

\subsection{Stage 4: time curriculum}
Extend the physics time interval through current-sheet formation and the first reconnection event. Physics difficulty is not uniform in time.

\subsection{Stage 5: full training-case family}
After one case is stable, train over all frozen training trajectories with case-balanced labeled and physics contexts.

\subsection{Stage 6: optional global and adaptive components}
Only after the local residual is verified should global balances, residual-adaptive sampling, or adaptive loss weights be added.

\subsection{Optimizer and checkpoints}
Adam is the primary optimizer for consistency with Module~4. Large full-field FNO objectives make deterministic quasi-Newton refinement more memory-intensive than in small coordinate PINNs, so it is not required in the primary specification. Any secondary optimizer must be reported separately.

Checkpoint selection uses a predeclared validation score that combines field accuracy with independent physics diagnostics. The final test set remains untouched.

\section{Gradient diagnostics and optimization pathology}
\label{sec:gradients}

\subsection{Per-term gradient norms}
For active component $\mathcal{L}_i$, define
\begin{equation}
g_i=\left\|\nabla_\theta\mathcal{L}_i\right\|_2.
\end{equation}
Track at least $g_{\rm data}$, $g_{\rm PDE}$, $g_{\rm IC}$, and $g_{\rm gauge}$ at representative intervals.

\subsection{Gradient alignment}
For losses $i$ and $j$,
\begin{equation}
c_{ij}=
\frac{\nabla_\theta\mathcal{L}_i\cdot\nabla_\theta\mathcal{L}_j}
{\|\nabla_\theta\mathcal{L}_i\|_2\,\|\nabla_\theta\mathcal{L}_j\|_2}.
\end{equation}
A negative value means the local descent directions conflict. Persistent data-PDE conflict can indicate poor scaling, limited capacity, reference-model discrepancy, or an optimization path problem \cite{Wang2021Gradients}.

\subsection{Layerwise gradients}
Inspect gradients in the lifting layer, spectral weights, local branches, and projection head. A physics loss that only modifies the final projection weakly may fail to reshape the latent operator representation.

\subsection{Spectral gradient diagnostics}
Because the FNO has explicit retained Fourier weights, it is useful to inspect whether physics gradients concentrate on low modes while current-sheet errors remain in unresolved high modes. Physics weighting cannot update spectral degrees of freedom that the architecture does not possess.

\subsection{Loss stiffness}
The vorticity residual combines advection, magnetic forcing, and a fourth-order viscous contribution. Their scales can differ sharply as current sheets form, producing a stiff optimization landscape even when the physical PDE is well posed.

\section{What physics information can and cannot do}
\label{sec:capability}

\subsection{Physics narrows the admissible operator family}
Data-only training asks the FNO to match observed input-output pairs. Physics-informed training additionally prefers operators whose outputs satisfy reduced-MHD dynamics between and around those labels.

\subsection{Physics contexts are not new truth arrays}
A residual at an unlabeled time says whether a prediction satisfies the governing equation. It does not reveal the exact field that the reference solver would produce there. Equation information and observation information are different resources.

\subsection{Physics can constrain interpolation between saved times}
If solver outputs are sparse in time, the PINO can sample intermediate $t_f$ values and penalize equation defects there. This is a plausible source of data efficiency. It must be demonstrated by learning curves rather than assumed.

\subsection{Physics does not create arbitrary initial-state diversity}
The PDE is valid for many initial conditions, but the FNO's learned dependence on the input field still comes from the training process. Version~1 exposes only the fixed $\sin x\sin y$ magnetic shape and the amplitude-scaled $\cos x-\cos y$ flow shape. Strong operator claims require Module~2 to generate additional verified initial-state families.

\subsection{Physics cannot repair an inadequate spectral bottleneck automatically}
If the network discards spatial modes needed to represent a current sheet, no loss weight can recover a degree of freedom absent from the model class. Architecture adequacy remains a prerequisite.

\subsection{Model discrepancy becomes central with real M3D-C1 data}
If future data are generated by equations containing physics absent from the reduced model, a strong reduced-MHD residual may actively pull predictions away from valid higher-fidelity data. Physics-informed learning is only as appropriate as the physical model imposed.

\section{Failure modes and signatures}
\label{sec:failure}

\subsection{Wrong sign convention}
A sign error in $J$, $\omega$, a bracket, or a diffusion term can produce apparently smooth training while enforcing the wrong dynamics. Analytic sign tests are mandatory.

\subsection{Broken time-scaling chain rule}
If $t$ is standardized but Equation~\eqref{eq:time_chain} is not respected, the time derivative is multiplied by the wrong factor. The operator can then compensate by distorting spatial terms.

\subsection{FFT wave-number or axis mismatch}
Swapped axes, incorrect negative-frequency ordering, or a mismatch between tensor layout and FFT dimensions corrupts all spatial derivatives while leaving output images visually plausible.

\subsection{Residual aliasing}
Nonlinear bracket products can fold unresolved frequencies into lower modes if dealiasing is inconsistent with Module~2. A spurious low residual can result from canceling numerical artifacts rather than physical terms.

\subsection{Insufficient FNO retained modes}
The data loss may be low while $J$, $\omega$, and residuals remain poor because high spatial frequencies are absent from the learned spectral branch.

\subsection{Loss-term domination}
One objective can overwhelm the others through gradient magnitude even when its scalar value looks comparable. Per-term gradient logging is required.

\subsection{Physics-context undercoverage}
A model can overfit a fixed set of physics times. Keep a fixed independent validation physics set and inspect dense time-resolved residuals.

\subsection{Gauge drift}
The PDE residual can remain small while the spatial mean of $U$ drifts because the local equations do not determine that constant.

\subsection{Initial identity failure}
A time-conditioned operator that does not reproduce its supplied initial state at $t=0$ is not representing the intended solution operator even if later data fit looks good.

\subsection{Unlabeled test-case leakage}
Using held-out test initial fields in a physics loss makes the training transductive. This can artificially improve ``unseen-case'' performance.

\subsection{Residual improvement with worse field error}
A PINO can move away from discrete solver targets toward a lower continuous/spectral residual, especially when reference discretization error is nonzero. Report both field error and residual instead of assuming they move together.

\subsection{Trivial low-amplitude dynamics}
An overly strong regularization or physics weighting can bias predictions toward smoother, lower-amplitude states that reduce some residual components while missing the actual nonlinear event.

\subsection{False time-extrapolation confidence}
A differentiable time input lets the operator return a number for any $t$. It does not guarantee physically meaningful prediction beyond the training horizon.

\section{Diagnostics during training}
\label{sec:diagnostics}

At minimum, log the following on training and fixed validation sets:
\begin{itemize}
  \item $\Ldata$, $\Lpde$, $\Lic$, $\Lgauge$, and active global losses;
  \item separate $R_\psi$ and $R_\omega$ RMS and maximum values;
  \item per-term gradient norms and selected gradient cosines;
  \item relative $L_2$ errors of $\psi$ and $U$ on labeled validation fields;
  \item derived $J$ and $\omega$ errors;
  \item magnetic and kinetic energy errors;
  \item peak $|J|$, current-sheet width, and peak-time error where defined;
  \item reconnection-flux and reconnection-rate errors;
  \item Fourier spectral tails of predicted fields and residuals;
  \item mean $U$ and mean $\psi$ defects;
  \item optimizer learning rate, gradient clipping events, and nonfinite counts;
  \item GPU/CPU memory, wall-clock time, and FFT/AD evaluation counts; and
  \item checkpoint identity and random seed.
\end{itemize}

A training dashboard should not show only one total scalar loss. The scientific question is multi-dimensional.

\section{Controlled experiments and ablations}
\label{sec:experiments}

\subsection{Matched Module 4 versus Module 5}
The primary comparison holds fixed:
\begin{itemize}
  \item complete train/validation/test trajectory manifests;
  \item labeled operator pairs and their budgets;
  \item preprocessing transforms;
  \item FNO input/output tensor contract;
  \item width, layer count, retained modes, activation, and padding rule;
  \item data-loss definition;
  \item primary evaluation metrics; and
  \item validation protocol.
\end{itemize}
Module~5 adds physics contexts, physics losses, and their compute cost.

\subsection{Data-budget curves}
Vary the number of labeled training trajectories and/or labeled target times per trajectory while keeping the physics protocol declared. Plot held-out field and physical errors versus labeled reference cost.

\subsection{Physics-context-budget curves}
At fixed labeled data, vary the number of unlabeled physics times per training case. Report both accuracy and runtime because full-field residual evaluation is expensive.

\subsection{Warm start versus from scratch}
Compare physics refinement of the frozen Module~4 checkpoint with random-initialized composite training. Count pretraining compute explicitly.

\subsection{Flux-only versus full two-equation physics}
The flux equation is cheaper and lower order. Compare
\begin{equation}
\Ldata+\mathcal{L}_{R_\psi}
\end{equation}
with the full two-equation objective to quantify the value and cost of the vorticity branch.

\subsection{Uniform versus adaptive physics times}
Compare uniform time sampling with residual-adaptive or event-stratified physics-context sampling under the same context budget.

\subsection{Global-loss ablation}
Add energy and mean-square-flux constraints only after the local residual baseline is stable. Report whether they improve held-out diagnostics or merely reduce their own scalar defects.

\subsection{Interpolation versus extrapolation}
Separate tests for coefficient interpolation, coefficient extrapolation, time extrapolation, held-out $A_0$ amplitudes, and genuinely new initial-field shapes. These are not interchangeable.

\subsection{Statistical reporting}
Train multiple seeds. Report mean, spread, and failed runs. Physics-informed objectives often increase optimization variability, so best-seed reporting is especially misleading.

\section{Worked analytic example: resistive decay as an operator test}
\label{sec:example}

\subsection{Exact solution}
Set $U_0=0$ and
\begin{equation}
\psi_0(x,y)=\sin x\sin y.
\end{equation}
Because $\Lap\psi_0=-2\psi_0$ and all brackets vanish, the exact finite-resistivity solution is
\begin{equation}
\psi(x,y,t)=e^{-2\eta t}\sin x\sin y,
\qquad U(x,y,t)=0.
\label{eq:decay_exact}
\end{equation}
The exact operator therefore maps
\begin{equation}
(\psi_0,0,\eta,\nu,t)
\longmapsto
(e^{-2\eta t}\psi_0,0).
\end{equation}

\subsection{Residual of a wrong decay rate}
Suppose the operator predicts
\begin{equation}
\widehat\psi=e^{-\alpha t}\sin x\sin y,
\qquad \widehat U=0.
\end{equation}
Then
\begin{equation}
\widehat\psi_t=-\alpha\widehat\psi,
\qquad
\Lap\widehat\psi=-2\widehat\psi.
\end{equation}
Therefore
\begin{equation}
\boxed{
\Rpsi=(2\eta-\alpha)\widehat\psi.
}
\label{eq:wrong_decay}
\end{equation}
The residual detects the wrong time dependence even if one snapshot happens to match the reference closely.

\subsection{Spectral implementation check}
The field contains only modes $(\pm1,\pm1)$. Fourier differentiation should reproduce the factor $-2$ to roundoff on the accepted grid. This example simultaneously tests time AD, Fourier Laplacian signs, physical $\eta$, and residual assembly.

\subsection{Gauge-invisible error}
If the operator predicts
\begin{equation}
\widehat U=C(t)
\end{equation}
with spatially constant $C$, then velocity, vorticity, and local PDE residuals are unchanged. The gauge loss is therefore necessary to detect this error.

\subsection{Small field ripple, large current error}
Let
\begin{equation}
\widehat\psi=\psi+\epsilon\sin(kx).
\end{equation}
The field perturbation is $O(\epsilon)$, but the current perturbation is
\begin{equation}
\delta J=\epsilon k^2\sin(kx).
\end{equation}
High-wave-number errors can therefore dominate derivative and physics diagnostics while remaining visually negligible.

\section{Recommended version-1 specification}
\label{sec:specification}

\begin{conventionbox}[\,: primary Module 5 starting point]
\begin{itemize}
  \item \textbf{Map:} $(\psi_0,U_0,\eta,\nu,t)\mapsto(\widehat\psi_t,\widehat U_t)$ on complete numerical fields.
  \item \textbf{Architecture:} the matched Module~4 2D FNO; starting specification width 32, four FNO blocks, retained modes $16\times16$, GELU, two-channel projection, no nonperiodic padding.
  \item \textbf{Input channels:} $[\psi_0,U_0,\eta,\nu,t]$, with scalar context broadcast spatially after training-only preprocessing.
  \item \textbf{$A_0$:} retained as metadata and an initial-field consistency check; not duplicated as a primary input channel.
  \item \textbf{Outputs:} only $(\psi,U)$; derive $J=-\Lap\psi$ and $\omega=\Lap U$.
  \item \textbf{Data:} exactly the frozen Module~4 labeled pair manifest and $\Ldata$.
  \item \textbf{Spatial physics:} Module~2 FFT wave numbers, spectral derivatives, and rectangular two-thirds dealiased brackets.
  \item \textbf{Time physics:} AD with respect to raw nondimensional target time through the complete preprocessing/network/inverse-scaling wrapper.
  \item \textbf{PDE:} normalized $\Rpsi$ and $\Romega$ on complete predicted grids for case-balanced unlabeled physics contexts.
  \item \textbf{IC:} soft full-field identity loss first; hard operator ansatz only as an ablation.
  \item \textbf{Periodicity:} structural periodic-grid representation; $\Lbc=0$.
  \item \textbf{Gauge:} soft full-grid zero-mean $U$ loss.
  \item \textbf{Global physics:} off in the primary endpoint; energy, mean-flux, and mean-square-flux terms are controlled ablations.
  \item \textbf{Physics sampling:} training cases only; continuous or stratified times within the declared training horizon.
  \item \textbf{Initialization:} primary warm start from the selected Module~4 checkpoint; from-scratch PINO as a required ablation.
  \item \textbf{Optimizer:} Adam with validation-selected schedule; all additional pretraining and residual compute counted.
  \item \textbf{Reporting:} repeated seeds, complete loss/gradient histories, held-out full fields, derivative-sensitive physics diagnostics, and compute cost.
\end{itemize}
\end{conventionbox}

No universal learning rate, batch size, physics-context count, or loss weight is frozen in this physics note. Those quantities depend on the measured dataset scales, memory budget, and validation behavior. They must be selected through documented training-case experiments rather than invented in advance.

\section{Implementation sequence and acceptance tests}
\label{sec:acceptance}

\subsection{Implementation sequence}
\begin{enumerate}[label=\textbf{M5.\arabic*.}]
  \item Load Module~3 trajectory manifests, tensor transforms, and the frozen Module~4 labeled-pair lists.
  \item Load the selected Module~4 FNO architecture and checkpoint; verify identical data-only inference.
  \item Implement a raw-nondimensional operator wrapper that performs preprocessing and inverse output scaling inside the differentiable graph.
  \item Implement Module~2-consistent FFT wave numbers, Laplacian, first derivatives, and two-thirds projection.
  \item Implement $\widehat J$, $\widehat\omega$, dealiased brackets, $\Rpsi$, and $\Romega$ using named intermediate tensors.
  \item Implement AD with respect to raw scalar $t$ and verify $\widehat\omega_t=\Lap\widehat U_t$.
  \item Implement separate labeled-pair, physics-context, IC, gauge, and optional global loaders.
  \item Implement residual normalization and per-term gradient diagnostics.
  \item Pass analytic Fourier and manufactured-time tests before learned training.
  \item Train one case through the staged curriculum.
  \item Compare PINO residuals with the accepted reference residual floor.
  \item Train the full frozen training-case family and run the controlled ablation matrix.
  \item Freeze the final validation-selected configuration before evaluating test cases.
\end{enumerate}

\subsection{Mandatory analytic tests}
\begin{itemize}
  \item Fourier first derivatives and Laplacians match analytic trigonometric modes near numerical precision.
  \item $J=-\Lap\psi$ and $\omega=\Lap U$ produce the frozen Module~2 signs.
  \item Dealiased brackets satisfy antisymmetry and $[f,f]=0$ for resolved test modes.
  \item The exact decaying mode in Equation~\eqref{eq:decay_exact} produces negligible $\Rpsi$ and $\Romega$.
  \item The wrong-decay test reproduces Equation~\eqref{eq:wrong_decay}.
  \item Raw-time AD agrees with an analytic manufactured time dependence after preprocessing.
  \item A constant offset in $U$ leaves local PDE residuals unchanged and is detected by $\Lgauge$.
  \item Every active loss produces finite gradients in relevant lifting, spectral, local, and projection weights.
\end{itemize}

\subsection{Mandatory numerical tests}
\begin{itemize}
  \item Residuals are invariant, after correct inverse transformation, under an equivalent change of neural input/output scaling.
  \item Batched physics-context evaluation agrees with an unbatched calculation on the same contexts.
  \item Case balancing is unchanged when duplicate times from one case are inserted deliberately.
  \item The physics operator uses the same FFT axis order and two-thirds mask as Module~2.
  \item $t=0$ identity fields survive HDF5-to-tensor-to-network preprocessing without channel or axis swaps.
  \item Checkpoint reload reproduces predictions, residuals, loss components, and preprocessing state.
  \item Fixed validation physics contexts remain unchanged across epochs even when training times are resampled.
  \item Precision and memory benchmarks are recorded for time AD plus full-grid spectral residuals.
\end{itemize}

\subsection{Acceptance criteria before production Module 5 training}
The full training family may begin only when:
\begin{enumerate}[label=\arabic*.]
  \item Module~2 reference cases have passed solver verification and convergence checks;
  \item Module~3 has frozen the full-field schema, trajectory splits, tensor ordering, and transforms;
  \item Module~4 has frozen its competent data-only FNO configuration and labeled budgets;
  \item all analytic spectral and time-AD residual tests pass;
  \item one training case reaches stable finite data and physics losses on independent validation physics times;
  \item predicted spectra show no unexplained pile-up at the retained or dealiased cutoffs;
  \item IC, gauge, mean, and energy diagnostics are internally consistent;
  \item test trajectories have not been used for residual scales, physics times, weights, or stopping rules; and
  \item the complete configuration is reproducible from saved metadata and seeds.
\end{enumerate}

\subsection{Minimum run artifacts}
Each run directory stores:
\begin{itemize}
  \item architecture and retained-mode configuration;
  \item Module~4 source checkpoint identifier when warm-started;
  \item training/validation/test case manifests;
  \item labeled-pair and physics-context sampling rules;
  \item preprocessing and residual scales;
  \item all loss weights and any adaptation histories;
  \item optimizer, learning-rate, batch, precision, and seed settings;
  \item complete component-loss and gradient histories;
  \item field, residual, spectral, energy, current-sheet, and reconnection diagnostics;
  \item best and final checkpoints;
  \item wall-clock, memory, FFT count, and AD timing summaries; and
  \item software revision and dependency metadata.
\end{itemize}

\section{Common misconceptions}
\label{sec:misconceptions}

\subsection{``A full-field neural operator is automatically physics-informed''}
False. Module~4 is a full-field FNO but remains strictly data-only. Physics enters only when a physical constraint affects training or architecture.

\subsection{``The input fields are images''}
False. They are numerical scalar fields with physical coordinates, normalization, sign conventions, and spectral meaning. A rendered image is a visualization product.

\subsection{``Full-field input means arbitrary initial conditions are solved''}
False. Version~1 still spans a narrow analytic initial-condition family. Operator input format and operator generalization are different concepts.

\subsection{``The PDE residual now has only second derivatives''}
False. The vorticity equation still contains fourth spatial derivatives of $U$. Fourier differentiation changes how they are computed, not the mathematical order.

\subsection{``Automatic differentiation computes all derivatives''}
Not in the primary formulation. AD computes time derivatives and training gradients. Fourier operators compute spatial derivatives.

\subsection{``One physics context on a $128^2$ grid is 16,384 new data points''}
False. It supplies 16,384 spatial equation defects but no independent reference field values.

\subsection{``No labels means a held-out test case can be used for physics training''}
False for a strict unseen-case experiment. Using its initial field and parameters during training makes the setting transductive.

\subsection{``A low PDE loss proves correct reconnection''}
False. Reconnection requires correct derivatives, topology, event timing, and current-sheet structure on held-out fields.

\subsection{``Equal scalar losses have equal influence''}
False. Weight updates are controlled by parameter gradients, not by scalar magnitudes alone.

\subsection{``The energy law can replace the local PDE''}
False. Global integrals discard spatial information and cannot uniquely determine the field.

\subsection{``More physics weight always produces more physical predictions''}
False. Excessive weighting can damage data fit, push toward smooth low-amplitude states, or amplify optimization stiffness.

\subsection{``Using the same FFT as the solver turns the PINO into the solver''}
False. The FFT is only a derivative and convolution tool. The network still learns a global time-conditioned solution operator by optimization rather than RK4 time marching.

\subsection{``Time as an input guarantees forecasting''}
False. Prediction beyond the trained time interval is extrapolation and must be tested explicitly.

\subsection{``This is now an M3D-C1 surrogate''}
False. The imposed physics and the data remain those of the synthetic two-field periodic reduced-MHD system.

\section{Module 5 computational contract}
\label{sec:contract}

\subsection{Input-output contract}
The primary model receives
\begin{equation}
(\psi_0[N_x,N_y],U_0[N_x,N_y],\eta,\nu,t)
\end{equation}
and returns
\begin{equation}
(\widehat\psi_t[N_x,N_y],\widehat U_t[N_x,N_y]).
\end{equation}
All arrays are numerical fields on the declared half-open periodic grid.

\subsection{Architecture contract}
Use the validation-selected Module~4 FNO family unchanged for the primary comparison. Any architecture modification is a separate ablation.

\subsection{Physics contract}
Use
\begin{align}
J&=-\Lap\psi,\qquad \omega=\Lap U,\\
\Rpsi&=\psi_t+[U,\psi]-\eta\Lap\psi,\\
\Romega&=\omega_t+[U,\omega]-[J,\psi]-\nu\Lap\omega.
\end{align}
Spatial derivatives use the Module~2 Fourier convention and primary nonlinear brackets use the rectangular two-thirds projection. Time derivatives use AD with respect to raw nondimensional $t$.

\subsection{Data and split contract}
Module~5 uses the identical Module~4 labeled operator pairs and complete physical-case split. Physics contexts may use training cases only in the primary unseen-case experiment.

\subsection{Loss contract}
The primary endpoint is
\begin{equation}
\Ltotal=w_d\Ldata+w_f\Lpde+w_0\Lic+w_g\Lgauge,
\qquad \Lbc=0,
\end{equation}
with global terms reserved for declared ablations. Residual normalization is applied only after the complete physical equation is assembled.

\subsection{Evaluation contract}
Evaluate complete untouched test fields using $\psi$, $U$, $J$, $\omega$, energy, spectra, current-sheet metrics, topology/reconnection diagnostics, gauge and mean defects, and dense PDE residuals. Report per-case results before aggregates.

\subsection{Fairness contract}
A claim that physics helped must account for identical labeled data and baseline architecture, additional physics-context evaluations, warm-start pretraining, residual derivative cost, hyperparameter search effort, and repeated-seed variability.

\section{Summary and bridge}
\label{sec:summary}

Module~5 converts the Module~4 data-only full-field FNO into a physics-informed neural operator without changing the prediction question. The network still maps complete initial fields and scalar context to complete target fields:
\[
(\psi_0,U_0,\eta,\nu,t)\mapsto(\widehat\psi_t,\widehat U_t).
\]
The difference is that predicted fields now participate in the reduced-MHD equations during training.

The full-field architecture changes the most natural derivative strategy. Spatial derivatives are computed by verified periodic Fourier operators, including two-thirds-dealiased nonlinear brackets, while time derivatives are obtained by automatic differentiation with respect to the scalar target time. This preserves the physical fourth-order structure of the vorticity equation without requiring four nested coordinate-AD levels.

The primary scientific comparison remains controlled: Module~4 and Module~5 share the same labeled pairs, whole-case split, preprocessing, FNO architecture, and data loss. Module~5 adds PDE, initial-state, and gauge information and must count the additional compute. Global balances are secondary ablations rather than replacements for local dynamics.

A successful result would show that the stated reduced-MHD information improves a full-field surrogate of the synthetic periodic model under a declared data and compute regime. It would not establish generalization to arbitrary plasma states, because the present initial-condition family remains narrow, and it would not establish fidelity to genuine M3D-C1 calculations.

\appendix

\section{Compact spectral residual ledger}
\label{app:ledger}

\subsection{Fourier identities}
For $k^2=k_x^2+k_y^2$,
\begin{align}
\FFT[f_x]&=ik_x\widehat f,\\
\FFT[f_y]&=ik_y\widehat f,\\
\FFT[\Lap f]&=-k^2\widehat f,\\
\FFT[J]&=k^2\widehat\psi,\\
\FFT[\omega]&=-k^2\widehat U,\\
\FFT[\Lap\omega]&=k^4\widehat U.
\end{align}

\subsection{Dealiased bracket ledger}
\begin{equation}
\mathcal{B}_{2/3}(f,g)=
\IFFT\left[
\Ptwo\FFT(f_xg_y-f_yg_x)
\right].
\end{equation}
Then
\begin{align}
\Rpsi&=\psi_t+\mathcal{B}_{2/3}(U,\psi)-\eta\Lap\psi,\\
\Romega&=\Lap U_t+\mathcal{B}_{2/3}(U,\Lap U)
-\mathcal{B}_{2/3}(-\Lap\psi,\psi)-\nu\Lap(\Lap U).
\end{align}
The second line is algebraically equivalent to the $J$-$\omega$ form and is useful for implementation audits.

\subsection{Periodic integral identities}
For smooth periodic $f,g,h$,
\begin{align}
\int_\Omega[f,g] \,\dd A&=0,\\
\int_\Omega f[g,h] \,\dd A
&=\int_\Omega g[h,f] \,\dd A
=\int_\Omega h[f,g] \,\dd A,\\
\int_\Omega f\Lap g\,\dd A
&=-\int_\Omega \nabla f\cdot\nabla g\,\dd A.
\end{align}
These identities underlie the energy and mean-flux checks.

\section{Loss and sampling ledger}
\label{app:lossledger}

\begin{longtable}{p{0.20\textwidth}p{0.33\textwidth}p{0.37\textwidth}}
\toprule
Quantity & Information required & Primary meaning \\
\midrule
\endfirsthead
\toprule
Quantity & Information required & Primary meaning \\
\midrule
\endhead
$\Ldata$ & Labeled target $\psi_t,U_t$ arrays & Supervised full-field fit \\
$\Lpde$ & Input fields, coefficients, time; no target field & Local reduced-MHD equation defect \\
$\Lic$ & Supplied input fields at $t=0$ & Solution-operator identity condition \\
$\Lgauge$ & Predicted $U$ full field & Fix spatial constant in stream function \\
$\Lbc$ & None in primary architecture & Zero because periodicity is structural \\
$\Lglobal$ & Predicted full fields and time derivatives & Optional integrated physical balances \\
Labeled pair count & Saved target fields & Amount of solver label information \\
Physics-context count & Unlabeled case-time queries & Number of operator evaluations constrained by PDE \\
Residual node count & Physics contexts times $N_xN_y$ & Spatial equation evaluations, not labels \\
Trajectory count & Independent accepted cases & Primary measure of physical case diversity \\
\bottomrule
\end{longtable}

\section{Symbol table}
\label{app:symbols}

\small
\begin{longtable}{p{0.16\textwidth}p{0.54\textwidth}p{0.20\textwidth}}
\toprule
Symbol & Meaning & Status \\
\midrule
\endfirsthead
\toprule
Symbol & Meaning & Status \\
\midrule
\endhead
$x,y$ & Periodic Cartesian coordinates & dimensionless \\
$t,T$ & Target time and final modeled time & dimensionless \\
$L_x,L_y$ & Spatial periods & $2\pi$ baseline \\
$N_x,N_y$ & Grid points & integers \\
$\psi$ & Magnetic-flux function & primary field \\
$U$ & Flow stream function & primary field \\
$J$ & Current variable $-\Lap\psi$ & derived field \\
$\omega$ & Vorticity $\Lap U$ & derived field \\
$\eta$ & Magnetic diffusivity & case/context scalar \\
$\nu$ & Kinematic viscosity & case/context scalar \\
$A_0$ & Initial stream-function perturbation amplitude & metadata; encoded in $U_0$ \\
$\mathcal{G}_\theta$ & Trainable full-field neural operator & model \\
$\theta$ & Trainable FNO parameters & model parameters \\
$\widehat\psi,\widehat U$ & Predicted target fields & model outputs \\
$\FFT,\IFFT$ & Discrete Fourier transform and inverse & numerical operators \\
$k_x,k_y$ & Fourier wave numbers & numerical \\
$k^2$ & $k_x^2+k_y^2$ & numerical \\
$\Ptwo$ & Rectangular two-thirds spectral projection & dealiasing \\
$\mathcal{B}_{2/3}$ & Dealiased Poisson-bracket operator & physics operator \\
$\Rpsi$ & Magnetic-flux equation residual & physics field \\
$\Romega$ & Vorticity equation residual & physics field \\
$S_{R_\psi},S_{R_\omega}$ & Whole-equation residual scales & training-only constants \\
$\Ldata$ & Supervised full-field data loss & objective \\
$\Lpde$ & PDE residual loss & objective \\
$\Lic$ & Initial identity loss & objective \\
$\Lgauge$ & Zero-mean stream-function loss & objective \\
$\Lglobal$ & Optional global balance loss & objective \\
$w_d,w_f,w_0,w_g,w_G$ & Top-level loss weights & hyperparameters \\
$E_B,E_K,E$ & Magnetic, kinetic, and total energy & diagnostics/physics \\
$R_E$ & Instantaneous energy-balance residual & optional physics \\
$g_i$ & Gradient norm of loss component $i$ & diagnostic \\
$c_{ij}$ & Cosine between loss gradients $i,j$ & diagnostic \\
\bottomrule
\end{longtable}
\normalsize

\section{Acronym and terminology table}
\label{app:acronyms}

\begin{longtable}{p{0.24\textwidth}p{0.66\textwidth}}
\toprule
Term & Meaning \\
\midrule
\endfirsthead
\toprule
Term & Meaning \\
\midrule
\endhead
AD & Automatic differentiation of a computational graph \\
Adam & Adaptive Moment Estimation optimizer \\
BC & Boundary condition \\
Case-balanced & Averaging or sampling that gives controlled weight to independent physical trajectories \\
Collocation / physics context & Unlabeled initial-field, coefficient, and time input at which a physics residual is evaluated \\
Data efficiency & Accuracy achieved for a declared amount of labeled reference information \\
Dealiasing & Removal/projection of unresolved quadratic product modes that would fold into lower frequencies \\
FNO & Fourier neural operator \\
FFT & Fast Fourier transform \\
Gauge & Convention fixing a physically redundant degree of freedom, here the spatial constant in $U$ \\
Hybrid PINO & Neural operator trained with both field labels and physical constraints \\
IC & Initial condition \\
MHD & Magnetohydrodynamics \\
MSE & Mean squared error \\
Neural operator & Learned map between field/function inputs and field/function outputs \\
Operator pair & One full initial-state/context input paired with one full target-time field \\
PDE & Partial differential equation \\
PINN & Physics-informed neural network, historically used for coordinate-network formulations \\
PINO & Physics-informed neural operator; the term used in this module for the full-field physics-informed model \\
Physics context & Unlabeled operator input used to evaluate a physical residual \\
Poisson bracket & $[f,g]=f_xg_y-f_yg_x$ \\
Residual & Equation defect after all PDE terms are moved to one side \\
RMHD & Reduced magnetohydrodynamics; the exact reduction and signs must be stated \\
Spectral convolution & Learned Fourier-mode channel mixing used by the FNO \\
Strong residual & Pointwise/grid representation of the differential equation defect \\
Time-conditioned operator & Operator that receives target time and predicts that state directly from the initial field \\
Transductive physics training & Training that uses unlabeled information from a nominally held-out test instance; not strict unseen-case evaluation \\
Warm start & Initialization from a previously trained checkpoint, here the matched Module~4 FNO \\
\bottomrule
\end{longtable}

\section{Concept questions and answer sketches}
\label{app:questions}

\subsection*{Questions}
\begin{enumerate}[label=\textbf{Q\arabic*.}]
  \item What exact change turns Module~4 into Module~5?
  \item What is the Module~5 full-field operator map?
  \item Why is $A_0$ not a separate primary input channel?
  \item Why does $A_0$ remain metadata?
  \item What is the difference between a coordinate PINN and the current PINO?
  \item Which derivatives are computed by AD in the primary formulation?
  \item Which derivatives are computed spectrally?
  \item Does spectral differentiation reduce the mathematical order of the vorticity equation?
  \item Why use the same FFT wave numbers and dealiasing convention as Module~2?
  \item Write $\Rpsi$ and $\Romega$ with the project signs.
  \item Why must outputs be inverse-scaled before residual construction?
  \item Why can transformed $\eta$ and $\nu$ channels not be used directly as PDE coefficients?
  \item What is an unlabeled physics context?
  \item Why are residual grid nodes not equivalent to reference labels?
  \item Why can using a test case only in $\Lpde$ still violate unseen-case evaluation?
  \item Why does the primary architecture need no boundary seam loss?
  \item Why is the gauge loss still needed?
  \item What is the advantage of sampling physics times between stored snapshots?
  \item Why can uniform residual MSE underweight a current sheet?
  \item Why is warm-starting from Module~4 scientifically useful?
  \item What additional cost must be counted for a warm-started PINO?
  \item Why can more physics weight make the result worse?
  \item What does a negative data-PDE gradient cosine indicate?
  \item Why can low $\psi$ error coexist with high $J$ error?
  \item What does holding out one $A_0$ value actually test?
  \item What would be needed to test generalization to new initial-field shapes?
  \item What must remain fixed in the primary Module~4/5 comparison?
  \item Why does time-conditioned output not guarantee forecasting?
  \item What claim about M3D-C1 is justified after this module?
\end{enumerate}

\subsection*{Answer sketches}
\begin{enumerate}[label=\textbf{A\arabic*.}]
  \item The same full-field operator is trained with declared reduced-MHD residuals or physical constraints that affect its weights.
  \item $(\psi_0,U_0,\eta,\nu,t)\mapsto(\widehat\psi_t,\widehat U_t)$ on complete numerical arrays.
  \item Version~1 already encodes the perturbation amplitude everywhere in $U_0=A_0(\cos x-\cos y)$.
  \item It is part of the physical case definition and permits provenance and consistency checks.
  \item A coordinate PINN maps coordinates to point values and differentiates them with respect to coordinates; the PINO maps complete input fields to complete output fields.
  \item Time derivatives of the predicted fields, plus ordinary backpropagation gradients with respect to network weights.
  \item First spatial derivatives, Laplacians, current, vorticity, and higher spatial derivatives required by the residual.
  \item No. The term $\Lap\omega=\Lap(\Lap U)$ is still fourth order; Fourier multipliers simply evaluate it directly.
  \item To keep spatial physics numerically consistent with the verified periodic solver and avoid quadratic aliasing ambiguity.
  \item $R_\psi=\psi_t+[U,\psi]-\eta\Lap\psi$ and $R_\omega=\omega_t+[U,\omega]-[J,\psi]-\nu\Lap\omega$.
  \item Neural scaling changes derivative factors; the PDE must be assembled from physical nondimensional fields.
  \item Log or standardized channels are network features, whereas the PDE coefficients are the physical nondimensional diffusivities.
  \item An initial-field/coefficient/time input with no target field, used only to evaluate physics defects.
  \item They are equation evaluations derived from the prediction, not independent observations of the true solution.
  \item The model has still seen the held-out initial state and parameters, making the setting transductive rather than unseen-case.
  \item The field lives on a half-open periodic grid and the spectral representation is periodically extended by construction.
  \item A spatial constant in $U$ changes no velocity, vorticity, or local PDE residual.
  \item The equation supplies training information where no solver target snapshot is stored, which can improve temporal data efficiency.
  \item A thin sheet occupies little area even though it controls peak current and reconnection.
  \item It starts Module~5 from the exact matched data-only solution, making subsequent changes more directly attributable to added physics.
  \item The complete Module~4 pretraining cost plus all physics-refinement derivative and optimization cost.
  \item Loss gradients can conflict, the residual can be misscaled, and excessive physics weighting can pull the model away from valid data.
  \item The two local descent directions oppose each other, indicating optimization conflict or model/data inconsistency.
  \item $J=-\Lap\psi$ multiplies a Fourier error by $k^2$, strongly amplifying high-frequency discrepancies.
  \item Amplitude interpolation or extrapolation within one fixed flow-field shape, not a new morphology.
  \item Module~2 must generate additional independently verified initial-field families and Module~3 must place them into whole-case splits.
  \item Labeled pairs, whole-case splits, preprocessing, FNO architecture, data loss, evaluation metrics, and comparable tuning opportunity.
  \item Numerical evaluation at arbitrary $t$ is possible, but accuracy beyond the trained time interval is extrapolation and must be measured.
  \item Only a physics-informed full-field surrogate of the stated synthetic two-field reduced-MHD family; genuine M3D-C1 fidelity is not established.
\end{enumerate}

\clearpage
\begin{thebibliography}{99}

\bibitem{Li2021}
Z. Li, N. Kovachki, K. Azizzadenesheli, B. Liu, K. Bhattacharya,
A. Stuart, and A. Anandkumar,
``Fourier neural operator for parametric partial differential equations,''
\emph{Ninth International Conference on Learning Representations} (2021),
\url{https://openreview.net/forum?id=c8P9NQVtmnO}.

\bibitem{Kovachki2023}
N. Kovachki, Z. Li, B. Liu, K. Azizzadenesheli, K. Bhattacharya,
A. Stuart, and A. Anandkumar,
``Neural operator: Learning maps between function spaces with applications to PDEs,''
\emph{Journal of Machine Learning Research} \textbf{24}, 1--97 (2023).

\bibitem{Raissi2019}
M. Raissi, P. Perdikaris, and G. E. Karniadakis,
``Physics-informed neural networks: A deep learning framework for solving forward and inverse problems involving nonlinear partial differential equations,''
\emph{Journal of Computational Physics} \textbf{378}, 686--707 (2019),
\url{https://doi.org/10.1016/j.jcp.2018.10.045}.

\bibitem{Karniadakis2021}
G. E. Karniadakis, I. G. Kevrekidis, L. Lu, P. Perdikaris, S. Wang, and L. Yang,
``Physics-informed machine learning,''
\emph{Nature Reviews Physics} \textbf{3}, 422--440 (2021),
\url{https://doi.org/10.1038/s42254-021-00314-5}.

\bibitem{Baydin2018}
A. G. Baydin, B. A. Pearlmutter, A. A. Radul, and J. M. Siskind,
``Automatic differentiation in machine learning: a survey,''
\emph{Journal of Machine Learning Research} \textbf{18}, 1--43 (2018),
\url{https://jmlr.org/papers/v18/17-468.html}.

\bibitem{Wang2021Gradients}
S. Wang, Y. Teng, and P. Perdikaris,
``Understanding and mitigating gradient flow pathologies in physics-informed neural networks,''
\emph{SIAM Journal on Scientific Computing} \textbf{43}, A3055--A3081 (2021),
\url{https://doi.org/10.1137/20M1318043}.

\bibitem{Krishnapriyan2021}
A. S. Krishnapriyan, A. Gholami, S. Zhe, R. M. Kirby, and M. W. Mahoney,
``Characterizing possible failure modes in physics-informed neural networks,''
in \emph{Advances in Neural Information Processing Systems 34} (2021),
\url{https://proceedings.neurips.cc/paper/2021/hash/df438e5206f31600e6ae4af72f2725f1-Abstract.html}.

\bibitem{Kingma2015}
D. P. Kingma and J. Ba,
``Adam: A method for stochastic optimization,''
\emph{Third International Conference on Learning Representations} (2015),
\url{https://arxiv.org/abs/1412.6980}.

\bibitem{Goedbloed2004}
J. P. Goedbloed and S. Poedts,
\emph{Principles of Magnetohydrodynamics: With Applications to Laboratory and Astrophysical Plasmas},
Cambridge University Press (2004).

\end{thebibliography}

\end{document}
